\section{Experimentos e demonstração}

\subsection{Setup experimental}

No enquadramento metodológico adotado, os experimentos correspondem à demonstração e à avaliação inicial do artefato de design. Seu objetivo é verificar se o \textit{pipeline} proposto consegue construir uma base consultável, recuperar evidências documentais, gerar respostas fundamentadas e preservar rastreabilidade suficiente para inspeção posterior, no escopo delimitado da 56\textordfeminine{} Legislatura do Senado Federal.

O experimento utilizou uma base derivada do \textit{dataset} utilizado no projeto\footnote{\href{https://huggingface.co/datasets/fabriciosantana/discursos-senado-legislatura-56}{\textit{Dataset} público no Hugging Face}.}, contendo 15.729 registros do conjunto de origem, dos quais 15.726 correspondiam a discursos efetivamente escritos e 3 foram descartados por ausência do texto integral. A escolha desse corpus dialoga com estudos que analisam computacionalmente discursos do Senado no período 2007--2024, mas recorta a 56\textordfeminine{} Legislatura para construir e avaliar uma base RAG específica \parencite{blanco2025plenarioPalanqueEstudio,blanco2026letrasGraudas}. A partir desses registros, foram gerados 23.806 \textit{chunks}, organizados em 120 arquivos Markdown de ingestão. Esses números foram registrados nos artefatos públicos de construção da base\footnote{\href{https://github.com/fabriciosantana/mcdia/blob/main/05-iag/4-project/knowledge_openwebui/build_metadata.json}{Metadados de construção da base}; \href{https://github.com/fabriciosantana/mcdia/blob/main/05-iag/4-project/knowledge_openwebui/discursos_chunks.jsonl}{arquivo JSONL de chunks}.}, assegurando rastreabilidade quantitativa da base construída.

Na etapa de preparação da base, adotaram-se como parâmetros principais \texttt{max\_words = 850}, \texttt{overlap\_words = 150} e \texttt{chunks\_per\_file = 200}, buscando equilibrar granularidade e preservação de contexto. Além do texto segmentado, cada \textit{chunk} foi enriquecido com metadados como data, autor, partido, unidade federativa, casa legislativa, tipo de uso da palavra, resumo, indexação e URL do discurso original, de modo a favorecer recuperação qualificada e verificação das respostas. Já na etapa posterior de ingestão no Open WebUI, os arquivos \textit{Markdown} previamente segmentados foram processados com \textit{Markdown Header Text Splitter} habilitado, \textit{Chunk Size} = 6000, \textit{Chunk Overlap} = 500 e \textit{Chunk Min Size Target} = 2000, além do uso do modelo de embedding \texttt{text-embedding-3-small} e lote de \textit{embeddings} de tamanho 32.

O \textit{prompt} de RAG foi evoluído iterativamente para atender a requisitos centrais do problema institucional: privilegiar o contexto recuperado, proibir explicitamente a invenção de fatos, exigir declaração de insuficiência quando a informação solicitada não estivesse presente no contexto, distribuir citações ao longo da resposta e exigir uma seção final de referências com os metadados dos discursos utilizados. O \textit{prompt} configurado no modelo RAG do Open WebUI está disponível no repositório do projeto\footnote{\href{https://github.com/fabriciosantana/mcdia/blob/main/05-iag/4-project/eval/prompts/rag_prompt.md}{\textit{Prompt} do fluxo RAG}.} e foi utilizado tanto no uso interativo do sistema na própria plataforma quanto na bateria automatizada. Nas consultas automatizadas, o \textit{script} enviou a pergunta do usuário e anexou a coleção indexada, deixando que o Open WebUI aplicasse o \textit{template} RAG configurado no servidor e realizasse internamente a recuperação, a montagem do contexto e a formulação da resposta. Com isso, buscou-se aproximar a avaliação automatizada do fluxo de uso validado manualmente, preservando fidelidade ao contexto recuperado, controle de escopo, reconhecimento explícito de ambiguidade e estrutura de saída orientada à verificação das fontes. Tal configuração alinha-se à literatura que enfatiza verificabilidade, controle documental e avaliação explícita em sistemas RAG \parencite{gao2023ragSurvey,saadFalconEtAl2024ares,esEtAl2023ragas}.

Para a avaliação, desenvolveu-se um \textit{script} em \textit{Python} que automatizou consultas à \textit{knowledge base} do Open WebUI via API, gerando saídas em \texttt{.jsonl}, \texttt{.md} e \texttt{.csv}. O script envia cada pergunta ao endpoint de chat da plataforma e anexa a coleção indexada como \textit{file} do tipo \texttt{collection}; com isso, a recuperação e a montagem do contexto final ficam a cargo do próprio Open WebUI, e não de lógica externa implementada no \textit{script}. Os artefatos preservam a pergunta, a resposta, as fontes recuperadas, a configuração da rodada, os identificadores da \textit{knowledge base} e os \textit{hashes} dos principais arquivos usados na preparação da base. A pontuação foi atribuída em uma segunda etapa de avaliação por LLM como juiz, segundo a estratégia de \textit{LLM as a Judge}, aplicada sobre rubrica com cinco critérios (\textit{aderência}, \textit{precisão factual}, \textit{foco nas fontes}, \textit{síntese} e \textit{confiabilidade/alucinação}), cada um com escala de 0 a 2\footnote{\href{https://github.com/fabriciosantana/mcdia/blob/main/05-iag/4-project/eval/RUBRIC.md}{Rubrica de avaliação automatizada}.}. O script permite configurar separadamente o modelo usado como LLM juiz; nas rodadas principais analisadas, por simplificação experimental, optou-se por reutilizar o mesmo modelo nas etapas de geração e avaliação. Essa combinação entre LLM juiz e rubrica explícita é compatível com a literatura sobre avaliação automatizada de geração textual e de saídas de RAG \parencite{liuEtAl2025judgeRag,liuEtAl2023geval}. Nas execuções principais, o modelo usado para geração e avaliação foi \texttt{gpt-5-nano}, com recuperação sobre coleção indexada no Open WebUI, e os \textit{embeddings} \texttt{text-embedding-3-small} permaneceram registrados nos metadados dos trechos retornados.

Assim, a avaliação combina quatro camadas complementares, mas não equivalentes. A recuperação é observada por meio dos metadados dos trechos e arquivos retornados pelo Open WebUI; a geração é examinada nas respostas completas produzidas pelo fluxo RAG; a avaliação automatizada atribui notas às respostas por meio de LLM como juiz e rubrica explícita; e a inspeção humana/manual atua como verificação qualitativa amostral de casos fortes, cautelosos, limítrofes e divergentes. Essa separação é importante porque uma nota alta na resposta final não demonstra, isoladamente, que a recuperação foi ótima, nem que o julgamento automatizado substitui validação humana independente.

Do ponto de vista da recuperação, o \textit{script} de avaliação não fixa explicitamente um \textit{top-k} nem aplica reranqueamento manual: esses passos permanecem internalizados na plataforma. Ainda assim, os artefatos da rodada base permitem observar o comportamento efetivo do sistema\footnote{\href{https://github.com/fabriciosantana/mcdia/blob/main/05-iag/4-project/eval/results/rag_eval_20260416T172816Z.question_analysis.md}{Análise da rodada base por pergunta}.}. Nessa rodada, cada resposta retornou metadados associados a 9 a 15 trechos-fonte, com média de 12,1, distribuídos por 3 a 13 arquivos \textit{Markdown} de origem, com média de 8,6 arquivos por resposta. Isso sugere que o contexto final utilizado na geração resultou da agregação automática de múltiplos trechos recuperados semanticamente a partir da coleção indexada.

A tipologia das perguntas não foi tomada integralmente de um único \textit{benchmark}, mas foi inspirada por \textit{benchmarks} de avaliação de RAG, especialmente o RAGEval, que organiza consultas em famílias como pergunta factual, integração de informações, comparação numérica, sequência temporal, raciocínio em múltiplas etapas, sumarização e perguntas não respondíveis \parencite{zhuEtAl2025rageval}. A partir dessa base, o presente experimento adaptou as categorias ao domínio parlamentar e aos riscos específicos observados no \textit{corpus} de discursos do Senado. Assim, categorias como \textit{factual}, \textit{comparação}, \textit{síntese}, \textit{multietapas} e \textit{evidência negativa} dialogam diretamente com a literatura, enquanto rótulos como \textit{foco no autor}, \textit{desambiguação}, \textit{autoria cruzada}, \textit{controle de escopo} e \textit{estresse de recuperação} devem ser entendidos como extensões metodológicas propostas neste trabalho, criadas para tensionar problemas especialmente relevantes em consultas sobre autoria, mistura indevida de discursos, controle de escopo e robustez da recuperação.

\subsection{Resultados alcançados}

Os resultados podem ser descritos em três frentes complementares: integridade da base, inspeção manual de respostas e bateria automatizada de perguntas. 

Na frente operacional, a base foi carregada e manteve-se operacional nas consultas realizadas, em coerência com os metadados gerados na preparação: 15.729 registros de entrada, 15.726 discursos aproveitados, 3 descartes por ausência de texto útil e 23.806 \textit{chunks} persistidos para recuperação semântica. Esses números não demonstram qualidade da resposta por si só, mas indicam que o experimento partiu de uma base operacionalmente íntegra e rastreável.

Na inspeção manual, adotaram-se três estudos de caso exemplificativos, documentados no relatório público de resultados manuais\footnote{\href{https://github.com/fabriciosantana/mcdia/blob/main/05-iag/4-project/eval-openwebui/RESULTS.md}{Relatório público de resultados manuais}.}, escolhidos para representar um caso forte, um caso de cautela e um caso de limite. No primeiro, a resposta sobre os argumentos de Paulo Paim contra a capitalização da Previdência mostrou desempenho verificável consistente: o sistema recuperou múltiplos discursos do senador, sintetizou os principais pontos e apresentou referências explícitas ao final. No segundo, diante de uma pergunta de desambiguação de autoria sobre Confúcio Moura, o sistema adotou postura cautelosa, reconhecendo insuficiência de contexto e evitando atribuições indevidas. No terceiro, diante de uma pergunta multifocal que exigia reconstrução de um plano legislativo consolidado para reduzir desigualdade, o sistema preservou parte dessa cautela, mas já operou em nível de abstração mais alto, revelando o limite do fluxo quando a consulta exige consolidação ampla, articulação entre autores e fechamento programático a partir de evidências fragmentárias.

A avaliação automatizada foi consolidada a partir das três rodadas completas executadas após o alinhamento entre o fluxo automatizado e a configuração operacional do Open WebUI\footnote{\href{https://github.com/fabriciosantana/mcdia/blob/main/05-iag/4-project/eval/results/rag_eval_20260416T172816Z.csv}{CSV da rodada base}; \href{https://github.com/fabriciosantana/mcdia/blob/main/05-iag/4-project/eval/results/rag_eval_20260416T182240Z.csv}{CSV da segunda rodada}; \href{https://github.com/fabriciosantana/mcdia/blob/main/05-iag/4-project/eval/results/rag_eval_20260416T232921Z.csv}{CSV da terceira rodada}; \href{https://github.com/fabriciosantana/mcdia/blob/main/05-iag/4-project/eval/results/stability_summary_20260417T012001Z.md}{síntese de estabilidade}.}. Em cada rodada, 20 perguntas foram submetidas à \textit{knowledge base}, todas as respostas foram geradas com sucesso e, em seguida, avaliadas em uma segunda etapa de avaliação por LLM como juiz, com base na rubrica definida para o experimento. As médias observadas foram 9,15, 9,55 e 9,35 pontos em 10, resultando em média agregada de 9,35. Não houve falhas nas três execuções. No vocabulário da rubrica adotada, 17 respostas da primeira rodada ficaram na faixa ``pronta para uso'' (8--10) e 3 na faixa ``utilizável com supervisão'' (5--7); nas duas rodadas seguintes, 19 respostas ficaram na primeira faixa e 1 na segunda. Essas faixas indicam qualidade estimada no escopo do experimento, e não prontidão operacional ampla. O conjunto de 20 perguntas, suas categorias e as notas obtidas nas três rodadas está disponível nos artefatos públicos da avaliação\footnote{\href{https://github.com/fabriciosantana/mcdia/blob/main/05-iag/4-project/eval/discursos_questions.json}{Conjunto de perguntas}; \href{https://github.com/fabriciosantana/mcdia/blob/main/05-iag/4-project/eval/results/rag_eval_20260416T172816Z.question_analysis.md}{análise por pergunta da rodada base}.}. Esse material já distingue categorias inspiradas diretamente na literatura de \textit{benchmarks} RAG \parencite{zhuEtAl2025rageval} e categorias introduzidas como adaptação metodológica deste trabalho.

\begin{table}[t]
\centering
\footnotesize
\caption{Síntese entre objetivos, evidências e interpretação dos resultados}
\label{tab:sintese-objetivos-resultados}
\begin{tabular}{@{}>{\raggedright\arraybackslash}p{0.24\textwidth}>{\raggedright\arraybackslash}p{0.38\textwidth}>{\raggedright\arraybackslash}p{0.30\textwidth}@{}}
\textbf{Objetivo avaliado} & \textbf{Evidência empírica} & \textbf{Interpretação} \\
\hline
Construir base verificável & 15.726 discursos aproveitados, 23.806 chunks gerados e metadados legislativos preservados & A base construída mostrou-se operacional e rastreável para a avaliação realizada. \\
Validar consulta assistida & Três estudos de caso manuais representando caso forte, caso de cautela e caso de limite & O sistema apresenta melhor desempenho quando a pergunta é delimitada por autor, tema ou evidência documental. \\
Avaliar desempenho automatizado & Três rodadas completas, com médias 9,15, 9,55 e 9,35, sem falhas operacionais & O desempenho agregado foi alto nas métricas adotadas, embora com variação em itens mais complexos. \\
Identificar limites de uso & Oscilação em perguntas de precisão numérica, controle de escopo, multietapas e autoria cruzada & Consultas de maior abstração exigem supervisão humana e formulação cuidadosa. \\
Verificar robustez da avaliação por LLM & Análise manual amostral com convergência geral, mas divergência relevante em item de checagem de alucinação & O LLM como juiz é útil para escalar a avaliação, mas não substitui validação humana em casos críticos. \\
\hline
\end{tabular}
\end{table}

Rodadas complementares com outro modelo na etapa de avaliação por LLM foram usadas apenas como teste. Embora tenham completado a bateria, elas atribuíram nota máxima a todos os itens, o que sugeriu um padrão excessivamente permissivo de julgamento por parte do LLM juiz alternativo. Por essa razão, tais execuções não foram adotadas como evidência quantitativa principal, mas como alerta metodológico sobre a influência do modelo avaliador na estratégia de \textit{LLM as a Judge}.

Os resultados por família de perguntas e por tipo de exigência cognitiva ajudam a qualificar melhor o comportamento do sistema, embora esta bateria ainda seja pequena para caracterizar de modo estável todas as categorias consideradas. Nesta amostra de 20 perguntas, consultas das categorias \textit{factual}, \textit{proposta}, \textit{foco no autor}, \textit{desambiguação} e \textit{evidência negativa} mantiveram desempenho alto e pouca variação entre as rodadas, o que sugere boa aderência ao contexto quando a tarefa exigia localizar autoria, sintetizar argumentos delimitados ou reconhecer ausência de evidência suficiente. Em contrapartida, as perguntas que mais tensionaram o sistema foram aquelas que exigiam precisão numérica, controle de escopo, reconstrução multietapas, comparação entre autores e distinção fina entre formulações semelhantes. A maior variação ocorreu na pergunta \texttt{q15}\footnote{\href{https://github.com/fabriciosantana/mcdia/blob/main/05-iag/4-project/eval/discursos_questions.json}{Conjunto público de perguntas da bateria automatizada}.}, voltada à recuperação de informação numérica específica, cujas notas foram 6, 10 e 10; também apresentaram oscilação relevante a \texttt{q16}, que tensionava controle de escopo, com notas 7, 9 e 10, e a \texttt{q18}, formulada para testar autoria cruzada, com notas 8, 7 e 9. Esses casos sugerem que parte da instabilidade não decorre de falha operacional, mas da dificuldade intrínseca de avaliar respostas que combinam múltiplas fontes, números dispersos, inferências de autoria e limites de extrapolação.

Esses achados convergem com a avaliação qualitativa do projeto: perguntas mais delimitadas, ancoradas em um autor, tema ou argumento específico, tendem a produzir respostas mais fortes e verificáveis; perguntas amplas, multifocais ou formuladas de modo a induzir generalização aumentam o risco de mistura indevida de discursos, atribuições excessivas e perda de precisão factual. A análise manual amostral da rodada base, conduzida sobre sete respostas selecionadas, indicou convergência com o LLM juiz na maior parte dos casos, mas também revelou uma divergência importante em pergunta de checagem de alucinação, na qual a avaliação humana considerou a resposta inadequada por erro de escopo apesar da nota máxima atribuída automaticamente. Essa divergência não deve ser tratada apenas como ruído avaliativo: ela constitui achado metodológico relevante, pois mostra que o juiz automatizado pode superestimar respostas que parecem formalmente adequadas, mas violam o escopo documental esperado. Essa triangulação reforça que a avaliação por LLM como juiz é útil para escalar a análise, mas deve ser interpretada como instrumento auxiliar, especialmente em tarefas abertas ou metavaliativas. Em termos metodológicos, isso sugere que a utilidade da prova de conceito é mais clara para consulta assistida e exploração do acervo com apoio documental, mas que respostas de maior abstração ainda dependem de supervisão humana e de refinamentos adicionais no desenho de recuperação, no controle de escopo e na formulação do prompt \parencite{gao2023ragSurvey,wangEtAl2025segmentation,saadFalconEtAl2024ares}.

A inclusão de referências documentais ao longo da resposta e ao seu final mostrou-se viável por meio de instruções adequadas no \textit{prompt}, o que facilita a conferência posterior das informações. Em síntese, os resultados sustentam a viabilidade técnica inicial da prova de conceito e oferecem evidência mais concreta de onde o sistema apresenta utilidade e onde seus limites permanecem mais visíveis. Para tornar a demonstração menos abstrata, os exemplos completos de interação manual estão disponíveis em relatório público do projeto\footnote{\href{https://github.com/fabriciosantana/mcdia/blob/main/05-iag/4-project/eval-openwebui/RESULTS.md}{Relatório público de resultados manuais}.}, incluindo um caso de bom desempenho, um caso em que o sistema reconhece insuficiência de contexto e evita atribuições indevidas de autoria, e um caso mais desafiador em que a pergunta exige consolidação legislativa ampla sob forte restrição de contexto.

\subsection{Ameaças à validade}

Os resultados devem ser interpretados como evidência inicial de uma prova de conceito reprodutível, e não como \textit{benchmark} definitivo para sistemas RAG legislativos. A principal ameaça à validade interna decorre do tamanho reduzido da bateria automatizada, composta por 20 perguntas, e do balanceamento desigual entre categorias: algumas famílias aparecem com poucos exemplos, o que limita inferências estáveis por tipo de pergunta. Há também ameaça associada ao uso de \textit{LLM as a Judge}, pois os testes de sensibilidade indicaram que a pontuação pode variar de forma relevante conforme o modelo avaliador e a rubrica utilizada. A validade de construto é limitada pelo fato de a avaliação combinar notas automatizadas, inspeção manual e estudos de caso, sem ainda dispor de validação humana ampla e independente. Além disso, esta etapa não incluiu baseline lexical ou sem RAG, estudo de ablação dos metadados, métricas formais de recuperação nem protocolo independente de anotação humana; tais procedimentos ficam como prioridades para uma rodada experimental posterior. Por fim, a validade externa deve ser considerada com cautela, pois os resultados dependem da configuração específica do Open WebUI, do \textit{prompt} operacional, dos parâmetros de \textit{chunking}, da coleção indexada e do corpus de discursos da 56\textordfeminine{} Legislatura. Assim, os achados são mais bem compreendidos como demonstração reprodutível de viabilidade e limites em um domínio parlamentar específico, passível de replicação e expansão em estudos posteriores.

\subsection{Potencial de generalização}

Embora o experimento tenha sido desenhado para discursos da 56\textordfeminine{} Legislatura do Senado Federal, parte do pipeline tende a ser reaproveitável em outros acervos legislativos. São potencialmente transferíveis a preparação documental em lotes, o enriquecimento dos \textit{chunks} com metadados, a ingestão em base vetorial, o uso de \textit{prompt} orientado a evidências, o registro de artefatos de execução e a avaliação por LLM como juiz combinada à análise manual amostral. Por outro lado, permanecem específicos do \textit{corpus} analisado a estrutura dos campos de origem, os metadados disponíveis, a natureza dos pronunciamentos em plenário, a distribuição temporal dos discursos e a tipologia de perguntas construída para tensionar problemas de autoria, escopo e recuperação nesse domínio. Assim, a contribuição mais generalizável do trabalho não é um conjunto universal de resultados, mas um protocolo reprodutível e adaptável para construção e avaliação inicial de sistemas RAG em acervos legislativos.

A transferência do protocolo para outros acervos exigiria, portanto, nova caracterização documental e institucional. Diferenças de idioma, formato dos documentos, estabilidade das URLs, granularidade dos metadados, regime de dados abertos, atualização do corpus e infraestrutura disponível podem alterar tanto a recuperação quanto a verificabilidade das respostas. Em bases legislativas com retificações frequentes, documentos substituídos, discursos complementados posteriormente ou campos incompletos, a qualidade do sistema dependeria de regras explícitas de versionamento, descarte, atualização e auditoria da base.

\subsection{Impacto na administração pública}

O impacto potencial do \textit{chatbot} legislativo com RAG na administração pública pode ser analisado em pelo menos quatro dimensões. Em primeiro lugar, há potencial de contribuição para a transparência, na medida em que as respostas passam a ser ancoradas em discursos públicos e citáveis, facilitando a verificação por terceiros. Em segundo lugar, observa-se potencial de ganho de eficiência, uma vez que o sistema pode reduzir o tempo necessário para localizar informações relevantes em grandes acervos documentais. Em terceiro lugar, a solução pode contribuir para a memória institucional, ao facilitar a recuperação de posicionamentos anteriores, temas recorrentes e evidências documentais. Por fim, identifica-se potencial de apoio à decisão, na medida em que gabinetes e equipes técnicas podem dispor de uma camada adicional de mediação informacional para explorar o acervo legislativo de forma mais estruturada \parencite{martelloteViola2025organizacaoConhecimento,deAlmeidaSantos2025aiGovernance,wfd2023guidelinesAiParliaments}.

Por operar sobre fontes controladas e permitir referência explícita aos discursos da 56\textordfeminine{} Legislatura, a solução é mais compatível com exigências de \textit{accountability} do que o uso isolado de um \textit{chatbot} generalista. Essa característica assume especial relevância em contextos públicos, nos quais a confiabilidade da informação não pode ser dissociada da governança das fontes e da possibilidade de auditoria posterior.

Ainda assim, o potencial institucional identificado não equivale a prontidão operacional. Em eventual adoção por órgão público, seria necessário definir responsabilidades de curadoria da base, periodicidade de atualização, tratamento de discursos removidos ou retificados, trilhas de auditoria, política de retenção de logs, controle de versões do \textit{prompt} e mecanismos de contestação ou correção de respostas. Também seria necessário distinguir usos assistivos, voltados à exploração documental por servidores e pesquisadores, de usos voltados ao atendimento público, que exigiriam níveis mais altos de validação, supervisão humana e governança.
